\section{A worked example, by hand}
\label{sec:example}

To connect the formulae to the numbers the program prints, this section
computes one photometric case almost entirely by hand: a $V = 12.0$
(Vega) solar-type star, observed in $V$ through a 200-mm f/5 Newtonian
(secondary obstruction 70\,mm, optics efficiency 0.85) with a typical
CMOS camera (QE $\simeq 0.8$ at 5500\,\AA{}, gain setting giving
1.0\,e$^-$/ADU, read noise 1.5\,e$^-$, pixel 3.76\,\si{\micro\metre}),
under $2.5''$ seeing at airmass 1.2 from a suburban site
(SQM = 20.0, elevation 300\,m), for 60\,s in a $5''$-radius aperture.

\paragraph{1. Photon budget of the source.} A $V=0$ star delivers about
$1.0\times10^{4}$ photons\,s$^{-1}$\,cm$^{-2}$\,\AA$^{-1}$... more
precisely, Eq.~(\ref{eq:flam}) with $Z_\nu = 3640$\,Jy at 5500\,\AA{}
gives $F_\lambda = 3.63\times10^{-9}$\,erg\,s$^{-1}$cm$^{-2}$\AA$^{-1}$,
i.e.\ $F_\lambda/(hc/\lambda) \simeq 1005$
photons\,s$^{-1}$cm$^{-2}$\AA$^{-1}$. At $V=12$: $\times10^{-0.4\cdot12}
= 1.58\times10^{-5}$, so $1.59\times10^{-2}$
photons\,s$^{-1}$cm$^{-2}$\AA$^{-1}$.

\paragraph{2. Collecting area and effective bandwidth.} $A = \pi(20.0^2 -
7.0^2)/4 = 275.7$\,cm$^2$. The Bessell $V$ band has an effective width
of $\simeq 870$\,\AA{}. Through atmosphere ($T^X \simeq 0.87^{1.2}
\simeq 0.85$ at $V$), optics (0.85) and QE (0.80):
\[
\dot S_{\rm tot} \simeq 1.59\times10^{-2} \times 275.7 \times 870
\times 0.85 \times 0.85 \times 0.80 \;\simeq\; 2.2\times10^{3}\;
\mathrm{e^-\,s^{-1}} .
\]
The program, integrating the real solar-type spectrum against the real
$V$ curve, reports $\simeq 2.1\times10^{3}$\,e$^-$/s --- the hand
estimate is good to ten percent, which is the accuracy of using an
effective width.

\paragraph{3. Aperture and sky.} A $5''$ radius is $2\times$ the seeing
FWHM: a Moffat ($\beta=2.5$) profile puts 89\% of the light inside
(Sect.~\ref{sec:psf}), so $\dot S = 0.89 \times 2100 \simeq 1.9\times10^{3}$
e$^-$/s, and $S = 1.1\times10^{5}$\,e$^-$ in 60\,s. The aperture covers
$\pi\,25 = 78.5$\,arcsec$^2$. SQM 20.0 means the sky is
$m_{\rm sky} = 20.0$ mag/arcsec$^2$ in $V$, i.e.\ per arcsec$^2$ the sky
is a $20.0$-mag source: $\dot B = 2100 \times 10^{-0.4(20.0-12.0)}
\times 78.5 \simeq 2.6\times10^{2}$\,e$^-$/s, $B = 1.6\times10^{4}$
e$^-$ in 60\,s. With a plate scale of $0.776''$/pixel the aperture
holds $N_{\rm pix} \simeq 130$ pixels.

\paragraph{4. Noise budget.} Shot noise: $\sqrt{S + B} =
\sqrt{1.1\times10^5 + 1.6\times10^4} = 355$\,e$^-$. Read:
$\sqrt{130}\times1.5 = 17$\,e$^-$. Quantization:
$\sqrt{130}\times 1.0/\sqrt{12} = 3$\,e$^-$. Scintillation,
Eq.~(\ref{eq:scint}): $\sigma/I = 0.09\times20^{-2/3}\times1.2^{1.75}
\times e^{-300/8000}/\sqrt{120} = 1.5\times10^{-3}$, so
$\sigma_{\rm scint} = 1.5\times10^{-3}\times1.1\times10^{5} =
165$\,e$^-$. Total: $\sqrt{355^2 + 165^2 + 17^2 + 3^2} = 392$\,e$^-$.

\paragraph{5. Result.} $\mathrm{S/N} = 1.1\times10^{5}/392 \simeq 280$
--- a differential precision of $1086/280 \simeq 3.9$\,mmag per frame
against an equal comparison star ($\times\sqrt2$: 5.5\,mmag), fine for
a 25-mmag exoplanet transit at two-minute cadence. Note where the noise
lives: shot noise dominates, but scintillation is already 42\% of it at
$V=12$ on a 20-cm aperture --- brighten the star by three magnitudes
and scintillation rules alone, which is why the S/N curve of
Fig.~\ref{fig:snrbudget} flattens. The peak pixel holds
$\simeq 2100\times0.07\times60 \simeq 9\times10^{3}$\,e$^-$
(central-pixel fraction 0.07 for this sampling) plus sky and dark ---
far from a CMOS full well at low gain, so the 60-s frame is safe, as the
program's saturation columns confirm.

Every number above is printed by one run of the program; the exercise
shows there is no magic between the inputs and the outputs --- only the
chain of Sect.~\ref{sec:chain}, executed carefully.
